% from ega4-11-fr.tex
\oldpage[IV-3]{166}
canonique, il résulte de (8.2.11) qu'il existe un indice $\alpha$ et un ouvert
quasi-compact $U''_\alpha\subset X''_\alpha$ tels que
$U'=h_\alpha^{-1}(U''_\alpha)$~; en outre, en vertu de (8.5.2, (i)), on peut
supposer qu'il y a une section $t''_\alpha$ de $\mathcal F''_\alpha$ au-dessus
de $U''_\alpha$ telle que $t'$ soit l'image canonique de $t''_\alpha$. Quitte à
remplacer $S$, $X$, $Z_\lambda$, $f_\lambda$, $\mathcal F$,
$\mathcal G_\lambda$ et $u_\lambda$ par $S''_\alpha$, $U''_\alpha$,
$(f''_{\alpha\lambda})^{-1}(U''_\alpha)$ et les restrictions correspondantes
de $\mathcal F''_\alpha$, $\mathcal G''_{\alpha\lambda}$ et
$u''_{\alpha\lambda}$, on peut donc supposer que $U'=X'$, que $t'$ est
l'image canonique d'une section $t$ de $\mathcal F$ au-dessus de $X$ et que
l'homomorphisme $A\to A'$ est injectif, ou encore, si $p:S'\to S$ est le
morphisme correspondant, que l'homomorphisme
$\mathcal O_S\to p_*(\mathcal O_{S'})$ intervenant dans la définition de $p$
est injectif. Il en résulte aussitôt en vertu de la platitude de
$\mathcal G_\lambda$ sur $S$ (et en se ramenant par exemple au cas où
$Z_\lambda$ est affine sur $S$ (\textbf I, 1.6.3 et 1.6.5)) que
l'homomorphisme canonique
$\rho:\mathcal G_\lambda\to(g'_\lambda)_*(\mathcal G'_\lambda)$ est lui aussi
injectif. Mais l'homomorphisme composé
\[
\Gamma(X,\mathcal F)\xrightarrow{\Gamma(u_\lambda)}
\Gamma(Z_\lambda,\mathcal G_\lambda)\xrightarrow{\Gamma(\rho)}
\Gamma(Z'_\lambda,\mathcal G'_\lambda)
\]
est égal par définition à l'homomorphisme composé
\[
\Gamma(X,\mathcal F)\xrightarrow{\Gamma(\rho)}
\Gamma(X',\mathcal F')\xrightarrow{\Gamma(u'_\lambda)}
\Gamma(Z'_\lambda,\mathcal G'_\lambda)
\]
donc l'image de $t$ par ces homomorphismes composés est
$u'_\lambda(t')=0$~; en vertu de l'injectivité de l'homomorphisme
$\Gamma(Z_\lambda,\mathcal G_\lambda)\xrightarrow{\Gamma(\rho)}
\Gamma(Z'_\lambda,\mathcal G'_\lambda)$ on en conclut que $u_\lambda(t)=0$
pour tout $\lambda\in L$, d'où $t=0$ par hypothèse, et finalement $t'=0$.
\end{proof}

\begin{proposition}[11.9.13]
\label{IV.11.9.13-fr}
Les notations étant celles de (11.9.1) et (11.9.5), supposons que $X$ soit un
préschéma sur un corps $k$, et que, en posant $S=\operatorname{Spec}(k)$, on
ait $X'=X\times_S S'$, où $S'$ est un $k$-préschéma quelconque. Alors, si la
famille $(u_\lambda)_{\lambda\in L}$ est séparante, il en est de même de
$(u'_\lambda)$.
\end{proposition}

\begin{proof}
On peut se borner au cas où $S'=\operatorname{Spec}(A')$ est affine. Si $A'$
est une $k$-algèbre de type fini, le morphisme $g:X'\to X$ est plat et de
présentation finie, et on est donc dans les conditions d'application de
(11.9.10, (ii), b)). Dans le cas général, on considère $A'$ comme limite
inductive de ses $k$-sous-algèbres $A'_\alpha$ de type fini, et on applique à
chaque $A'_\alpha$ le résultat de (11.9.10, (ii), b))~; on conclut alors à
l'aide du lemme (11.9.12), puisque les $\mathcal G_\lambda$ sont $S$-plats.
\end{proof}

\begin{env}[11.9.14]
\label{IV.11.9.14-fr}
Gardons toujours les notations de (11.9.1) et (11.9.5) et supposons que $X$
soit un $S$-préschéma. Si pour tout changement de base
$g:X\times_S S'\to X$, où $S'\to S$ est un morphisme quelconque, la famille
$(u'_\lambda)$ correspondante est séparante, nous dirons que la famille
$(u_\lambda)$ est \emph{universellement séparante relativement à} $S$.
Lorsque la famille $(u_\lambda)$ est réduite à un seul élément $u$, nous
dirons encore que $u$ est \emph{universellement injectif}, relativement à
$S$. Il est clair alors que pour tout morphisme $h:S'\to S$, la famille
$(u'_\lambda)$ correspondante est universellement séparante relativement à
$S'$~; inversement, si $h$ est \emph{fidèlement plat} et si $(u'_\lambda)$ est
universellement séparante relativement à $S'$, alors $(u_\lambda)$ est
universellement séparante relativement à $S$, comme il résulte aussitôt de
(11.9.10, (i)) et du fait que pour tout morphisme $S''\to S$, le morphisme
correspondant $S''\times_S S'\to S''$ est fidèlement plat.
\end{env}

